% FICHIER GENERE PAR construire_v3.py, NE PAS EDITER A LA MAIN.
% Source : chapitres/13_inventaire_hypotheses.tex, moins les sections
% retirees, qui restent integralement dans main_v2. Toute correction
% se fait dans le fichier source puis par relance du script, sinon
% les deux divergent.
% ---------------------------------------------------------------------
% STATUT : RÉDIGÉ (transplanté de exploratory/cascade_qualitative/
%          chapitre_robustesse.tex, préambule retiré).
% FAIT (août 2026) : la ligne u_ij est au tableau des paramètres, avec sa
%          sensibilité mesurée par l'échelle de repli du ch. 10, ainsi qu'une
%          ligne pour l'élasticité sévérité/taille calibrée au ch. 5.
% ---------------------------------------------------------------------
\chapter{Robustesse et incertitude du chiffre reporté}
\label{ch:robustesse}

\noindent Le chapitre~\ref{chap:resultats} a conduit les tests de robustesse eux-mêmes
(tornado, sensibilités, invariance de l'ordre de remédiation). Ce chapitre ne les refait pas.
Il en donne la synthèse en une table, puis pose la question que ces tests laissent ouverte :
une fois l'incertitude reconnue, \emph{quel} nombre reporter. L'inventaire exhaustif des
paramètres, de leur nature et de leur sensibilité, est reporté au \S\ref{sec:ann-parametres}
pour qu'aucun chiffre du mémoire ne soit sans origine déclarée.

\begin{cle}
\textbf{Thèse.} Les conclusions du modèle ne reposent pas sur les paramètres non calibrés.
Je le montre en deux temps : (i)~elles sont insensibles aux choix de modélisation
innocents (structure de la cascade, famille de dépendance, couche temporelle, source de
sévérité) ; (ii)~elles ne réagissent qu'aux leviers qui doivent compter (conformité,
contagion), ce qui est une \emph{qualité}, pas un défaut, c'est la sensibilité au risque que
la Formule Standard n'a pas. Reste une incertitude irréductible, la queue de sévérité, que
j'assume en présentant le SCR comme une distribution, jamais un point.
\end{cle}

\section{L'argument de robustesse, en une table}

\begin{center}\footnotesize
\renewcommand{\arraystretch}{1.3}
\begin{tabular}{@{}p{4.5cm} p{4.6cm} p{4.4cm} c@{}}
\toprule
\textbf{Conclusion} & \textbf{Ce qu'on fait varier} & \textbf{Effet mesuré} & \textbf{Verdict} \\
\midrule
La non-conformité coûte cher ($\Delta_{\text{DORA}}>0$, grand) &
source de sévérité (OpRisk / PRC), famille de copule ($\theta\in[1,4]$) &
facteur ${\sim}1{,}9$ entre sources ; $\Delta_{\text{DORA}}$ varie $<1{,}4\,\%$ sur $\theta$ &
\rob \\
Priorité de remédiation P1 puis P4 (les \emph{sources}) &
cadre de dépendance (cascade vs copule), branchement vs marche &
P1 et P4 en tête dans tous les cadres ; SCR $\pm 4\,\%$ (branchement) &
\rob \\
C'est l'\textbf{asymétrie} de $W$ qui porte le résultat &
knockout de l'asymétrie (3\,000 tirages), placebo directionnel &
sans la direction, le signal tombe sous le bruit du placebo &
\rob \\
Pas de couche de Hawkes &
noyau exponentiel puis à retard, logique two-phase &
excitation intra-journalière (pic $3{,}3$ h) ; endogénéité de $0{,}551$ à ${\sim}0$ si les
co-occurrences du jour passent en exogène &
\rob \\
Les deux lois du socle prédisent une période qu'elles n'ont pas vue &
origine glissante, douze années notées hors échantillon &
fréquence : couverture $91{,}7\,\%$ pour la NB contre $75{,}0\,\%$ pour le Poisson, log-score en
sa faveur ; forme de queue non rejetée ($p=0{,}4087$) &
\rob \\
Le SCR \emph{réagit} à la conformité et à la contagion &
score de conformité $q\in[0,1]$, gain $g\in[0{,}05\,;0{,}90]$ &
$-53\,\%$ de NC à C ; $\times 1{,}6$ sur la plage de $g$ &
\feat \\
Le SCR est une distribution large, pas un point &
incertitude de paramètre (bootstrap deux niveaux) &
facteur ${\sim}2{,}5$ sur la VaR par la seule incertitude de calibration &
\ass \\
Le \emph{niveau} de sévérité ne tient pas hors échantillon &
notation d'une seule année à la fois, seuil ré-estimé sur l'apprentissage &
quantiles dépassés trois fois trop souvent ; l'échelle de queue dérive de $4{,}00\,\%$ par an,
soit au moins $+24{,}7\,\%$ sur le quantile &
\ass \\
Le \emph{signe} de $\Delta_{\text{DORA}}$ ne dépend pas de l'additivité des coûts d'un sinistre &
exposant $\theta\in[0{,}70\,;1{,}30]$ sur le nombre de piliers touchés &
écart de $9\,706$ à $20\,487$~M\euro{} : signe et ordre conservés, \emph{amplitude} non
($76\,\%$ de l'écart publié) &
\ass \\
\bottomrule
\end{tabular}
\end{center}

\vspace{4pt}
\textbf{Lecture.} Les cinq premières lignes sont l'ossature : le signe et l'ordre
de grandeur de l'effet DORA, la hiérarchie des piliers, le rôle de la direction, le statut de
l'auto-excitation et le comportement prédictif des deux lois du socle tiennent quand on change ce
qui pourrait sembler arbitraire. La sixième ligne
n'est pas une fragilité mais le c\oe ur du modèle : un SCR qui ne bougerait pas avec la
conformité serait inutile (c'est le reproche fait à la Formule Standard,
chapitre~\ref{chap:resultats}). Les trois dernières sont les limites honnêtes, et elles ne portent
pas sur le même objet : la queue de sévérité domine l'incertitude sur le \emph{niveau}, donc toute
lecture ponctuelle du SCR serait une surinterprétation ; la dérive de l'échelle de sévérité, seule
limite du mémoire découverte hors échantillon, déplace ce même niveau dans un sens connu et d'une
taille bornée (\S\ref{soc:sec:backtest}) ; l'additivité des coûts d'un sinistre
multi-piliers domine celle qui pèse sur l'\emph{écart}, sans jamais en renverser le signe
(\S\ref{sec:additivite-couts}).

% SOURCES-SCRIPTS: 25 29 47 63 61 66 67 74 08h 89



\section{Les trois étages de l'incertitude du SCR}
\label{sec:trois-etages}

L'incertitude sur le SCR ne se réduit pas au paramètre. Elle se répartit sur trois étages, qu'il
faut nommer plutôt que confondre (figure~\ref{fig:bande-modele}). Le premier est l'incertitude
de paramètre : l'intervalle de confiance de $\xi$ à lui seul dilate la VaR d'un facteur
$2{,}5$ sur la calibration figée, valeur que le bootstrap retrouve à la précision près.
% SOURCES-SCRIPTS: 63 46
Le deuxième est l'incertitude d'identification :
la direction de $W$ n'étant pas identifiable, la contagion se lit en bande
(chapitre~\ref{chap:partielle}). Le troisième, le plus souvent passé sous silence, est
l'incertitude de modèle : le choix de la \emph{famille} elle-même.

On la chiffre en recalculant le SCR sous des familles concurrentes. Sur l'axe de la
\textbf{dépendance}, à marges par pilier fixées, le SCR va de $5\,322$ M\euro{} (indépendance) à
$9\,806$ M\euro{} (comonotone, borne de Fréchet), soit un facteur $1{,}8$, la formule standard
Solvabilité~II et les copules gaussienne et de Gumbel se plaçant entre les deux ; notre cascade
dirigée ($8\,318$ M\euro) tombe dans cette bande, entre une dépendance gaussienne et une
dépendance de Gumbel. Sur l'axe de la famille de queue, l'écart est bien plus large :
d'une lognormale ($4\,330$ M\euro) à une GPD à queue lourde ($\xi=0{,}90$ ; $24\,016$ M\euro),
un facteur $5{,}5$. C'est le choix de la loi de queue qui domine, de loin, l'incertitude de
modèle, et c'est précisément pourquoi le plafond de réassurance ($40$ M\euro) est un levier
distinct, qui ramène la famille lourde à $495$ M\euro : une réponse au risque de queue, non le
risque lui-même.

\begin{figure}[htbp]\centering
\figover{J4_bande_modele.png}
\caption{La bande de modèle, troisième étage d'incertitude. (a)~axe dépendance, à
marges fixées : la cascade se place entre les copules gaussienne et de Gumbel, dans une bande
$[5{,}3\,;9{,}8]$ Md\euro. (b)~axe famille de queue : la queue lourde domine
l'incertitude ($\times 5{,}5$). (c)~les trois étages (paramètre, identification,
modèle) tels que le script les cumule ; la bande \emph{reportée} par le mémoire n'en retient que
les deux premiers, le troisième étant publié comme axe de sensibilité
(\S\ref{sec:etage-modele-hors-bande}).}
% SOURCES-SCRIPTS: 48
\label{fig:bande-modele}
\end{figure}

\subsection{Deux étages se cumulent dans la bande reportée, le troisième en sort}
\label{sec:etage-modele-hors-bande}

Les trois étages ne peuvent pas tous entrer dans un intervalle unique, et il faut dire lequel
en sort et pourquoi. La bande reportée par ce mémoire ne cumule que les deux premiers,
paramètre et identification. L'étage de modèle est publié, mais comme axe de
sensibilité déclaré et non comme composante de la bande.

\paragraph{Les trois étages, mis sur une seule échelle avant toute comparaison.} Comparer les
étages exige de les lire sur la \emph{même} grandeur, et c'est la première chose à faire ici :
l'étage de paramètre est publié sur le quantile d'\emph{un sinistre} (facteur $2{,}5$) tandis que
les deux autres vivent sur la \emph{charge annuelle agrégée}. Les rapprocher tels quels
comparerait deux échelles, ce qui est précisément l'erreur que ce mémoire s'interdit ailleurs.
Porté sur l'agrégat, l'intervalle bootstrap de $\xi$ donne $[3\,042\,;\,27\,152]$~M\euro, soit un
facteur $8{,}93$ : le quantile annuel amplifie l'incertitude de paramètre, et le facteur
$2{,}5$ ne mesure donc pas ce que le paramètre fait au \emph{capital}.
% SOURCES-SCRIPTS: 71

\begin{center}\footnotesize
\renewcommand{\arraystretch}{1.25}
\begin{tabular}{@{}l l r@{}}
\toprule
\textbf{Étage, lu sur la charge annuelle} & \textbf{Bande} & \textbf{Facteur} \\
\midrule
Paramètre ($\xi$ sur son IC$_{90}$) & $[3\,042\,;\,27\,152]$ & $8{,}93$ \\
Identification (direction de $W$) & $[6\,858\,;\,8\,697]$ & $1{,}27$ \\
Modèle, axe dépendance & $[5\,322\,;\,9\,806]$ & $1{,}84$ \\
Modèle, axe famille de queue & $[4\,330\,;\,39\,901]$ & $9{,}21$ \\
\bottomrule
\end{tabular}
\end{center}

\vspace{4pt}
\noindent En M\euro. La hiérarchie ainsi rétablie n'est pas celle qu'on annonce d'ordinaire :
\textbf{l'étage de modèle n'écrase pas les autres}, il est du même ordre que l'étage de paramètre,
et les deux dominent largement l'identification. Le retirer ne peut donc pas se justifier par sa
taille, et l'argument qui suit est d'une autre nature. Réserve à garder : ces largeurs de
paramètre font varier $\xi$ seul, $\sigma$ restant à sa valeur publiée, alors que le bootstrap
estime le couple conjointement ; ce sont donc des \emph{majorants}.

\paragraph{Ce qui justifie le déplacement, et ce n'est pas une taille.} Les deux étages retenus
sont de nature \emph{statistique} : ils mesurent ce que la donnée ne permet pas de trancher à
modèle fixé, et un intervalle est l'objet qui convient pour les porter. Le troisième est
\emph{épistémique} : il mesure l'effet d'un choix de famille, et un choix ne se moyenne pas.

Le point décisif est plus précis que cette opposition générale, et il se voit sur l'axe de la
famille de queue. Cet axe pousse $\xi$ à $0{,}90$, tandis que l'étage de paramètre le fait varier
jusqu'à $0{,}8313$ : c'est le même paramètre, et une bande doit dire \emph{une seule
fois} jusqu'où on le laisse aller. Or $0{,}8313$ est une borne estimée, celle de
l'intervalle bootstrap, quand $0{,}90$ est une valeur posée hors de cet intervalle. Les
placer sur la même échelle met un scénario et une borne de confiance dans le même objet, et le
lecteur ne sait plus si la borne haute est un aléa défavorable ou une autre théorie de la queue.
Le retrait rend la bande homogène, ce qui est un motif plus fort que la lisibilité.

À l'inverse, la borne \emph{basse} de cet axe, la lognormale à $4\,330$~M\euro, est une autre
famille et non un autre $\xi$ : celle-là est un apport propre, et c'est pourquoi le
recouvrement est asymétrique. Le déplacement porte donc sur ce que la bande \emph{annonce},
non sur ce que le mémoire publie.

\paragraph{Où cet étage est publié, désormais.} Il l'est deux fois, et rien n'est caché. Sur
l'axe de la famille de queue, la sixième posture de la table du
\S\ref{sec:predictive-robuste}, le pire cas de modèle à $\xi$ \emph{posé} à $0{,}90$, \emph{est}
l'ambiguïté de famille : c'est le même objet, sous un autre nom, et il porte un chiffre.
Sur l'axe de la dépendance, la bande de $5\,322$ à $9\,806$~M\euro{} reste publiée dans cette
section, avec le repère qui compte, à savoir que la cascade dirigée s'y place entre une
dépendance gaussienne et une dépendance de Gumbel.

\paragraph{Et ce qui ne change pas.} L'ordre d'amorce des piliers reste invariant à la
famille. Le déplacement porte sur la façon de \emph{reporter} le niveau, pas sur la thèse : le
niveau bouge dans la bande, l'ordre résiste, et c'est vrai des trois étages qu'on les cumule
ou non.

\paragraph{Sur l'auto-excitation, une convergence de choix et non une corroboration.} Le parti de
ne pas poser de couche de Hawkes est celui qui prête le plus le flanc, l'auto-excitation étant
devenue l'outil par défaut du cyber. Il est ici adossé à un test contre les variantes de noyau de
la littérature. Le préprint de cascade climatique cité au chapitre~\ref{ch:positionnement}
\citep{ccrn2026} arrive au même endroit, et \emph{pour la même raison de résolution temporelle} :
il retient un graphe acyclique à l'échelle de l'événement, au motif que le processus ponctuel
auto-excité devient naturel lorsque les instants exacts et la rétroaction récurrente sont au
c\oe ur du problème, ce qui n'est pas le cas quand la propagation se joue à l'intérieur d'un
événement dont on n'observe pas la chronologie fine. C'est exactement la situation du présent
travail, et la convergence est d'autant plus notable qu'elle est indépendante : autre péril,
autre équipe, et sans connaissance des travaux sur lesquels s'appuie le test conduit ici.

\begin{attention}
\textbf{Ce que cette convergence vaut, et ce qu'elle ne vaut pas.} Il serait commode de la lire
comme une corroboration du rejet ; ce serait la surinterpréter, et la précision importe puisque
c'est le point le plus exposé du mémoire. Ces auteurs ne testent pas l'auto-excitation contre
leur cascade : ils \emph{posent} le graphe acyclique, et leur mention du processus auto-excité
tient en une ligne d'un tableau de positionnement où il figure comme une représentation mature
parmi d'autres, mieux adaptée à un autre régime d'observation. Ce qui converge est donc le
\textbf{choix de modélisation et son motif}, pas une réfutation. Deux équipes qui ne disposent
pas de la chronologie fine d'un événement écartent le même outil pour la même raison, ce qui
déplace ce parti d'une préférence vers une contrainte d'observation, et c'est tout ce que l'on
peut en tirer. Le rejet lui-même ne repose que sur le test du présent travail.
\end{attention}

\paragraph{Ce que ce test établit réellement : une équivalence observationnelle, et c'est plus
fort qu'un rejet.} Le mot \og rejet \fg{} suppose que la donnée départage les deux modèles. Elle
ne les départage pas, et le mesurer valait mieux que de l'espérer.
% SOURCES-SCRIPTS: 85

L'argument se pose d'abord \emph{sans} rien simuler. Si chaque pilier touché est un événement
enregistrable, la fraction d'événements qui sont des \emph{enfants} vaut
$(\mathbb{E}[k]-1)/\mathbb{E}[k]$, où $\mathbb{E}[k]$ est le nombre moyen de piliers touchés par
sinistre : c'est exactement ce qu'un ajustement de Hawkes appelle son ratio de branchement. Avec
$\mathbb{E}[k]=1{,}931$ à l'état non conforme, une cascade sans aucune auto-excitation
prédit $n=0{,}482$, à comparer au $0{,}551$ mesuré sur la chronologie réelle.

Le test le confirme. On engendre soixante années de dates avec le moteur de cascade, dont la
seule structure temporelle est la durée des pas de propagation \emph{à l'intérieur} d'un sinistre,
puis on lui ajuste un Hawkes exponentiel comme on le ferait sur une chronologie réelle. À la
résolution de la donnée disponible, le jour, l'ajustement trouve $n=0{,}480$, soit la prédiction
analytique à trois millièmes. Il retrouve même l'horloge interne : la demi-vie estimée du noyau
vaut $9{,}0$~h pour un lag de propagation vrai de $6{,}7$~h, et il l'appelle \og décroissance de
l'excitation \fg{}. L'estimateur a été validé au préalable sur un Poisson pur, où il rend
$0{,}004$, et sur trois processus de Hawkes simulés de ratios connus, qu'il retrouve.

\begin{cle}
\textbf{Conséquence, à assumer telle quelle.} Un ratio de branchement de l'ordre de $0{,}5$
mesuré sur des dates réelles ne prouve pas l'auto-excitation : une cascade qui n'en
contient aucune le produit. Les deux modèles sont \emph{observationnellement équivalents} à la
résolution disponible. L'énoncé juste n'est donc pas \og le Hawkes est rejeté \fg{} mais \og le
Hawkes n'est pas distinguable de la cascade, et la cascade est retenue par parcimonie
interprétative \fg{} : elle rend compte du même regroupement temporel par un mécanisme qui porte
en outre la structure inter-piliers et se traduit en capital. C'est une conclusion plus faible en
apparence et plus solide en fait, puisqu'elle cesse de s'appuyer sur un départage que la donnée
ne permet pas.

\textbf{Attention au sens de l'effet de résolution.} Dégrader la résolution fait monter
le ratio estimé, de $0{,}473$ à l'heure jusqu'à $0{,}739$ au mois, parce qu'agréger fabrique du
regroupement apparent. C'est l'opération \emph{inverse} de celle conduite plus haut, qui fait
tomber l'endogénéité en reclassant les co-occurrences du même jour en exogène. Les deux
expériences ne se contredisent pas, elles déplacent deux choses différentes, et les confondre
donnerait une conclusion opposée.
\end{cle}

\begin{cle}
\textbf{Oui, chaque brique a des alternatives.} Fréquence (Poisson, Hawkes, Cox), sévérité
(lognormale, splicing, EVT bayésienne), dépendance (copules, réseaux, Hawkes multivarié),
mesure (TVaR, spectrale, robuste), agrégation (formule standard, Panjer) : on ne prétend pas
détenir \emph{la} méthode. La discipline est ailleurs : on a testé les concurrents les plus
sérieux, on borne l'effet du choix, et surtout le niveau bouge dans la bande mais
l'ordre d'amorce des piliers ne bouge pas ($P1>P4>P2>P3>P5$, invariant à la famille).
Le niveau est illustratif, l'ordre résiste : c'est, chiffrée sur un troisième axe, la thèse du
mémoire, et la réponse à qui demanderait \og et si vous aviez pris une autre méthode~?\fg.
\end{cle}

\section{Quelle grandeur reporter}
% PONT v3 : la section complete est dans main_v2, version
% longue et non modifiee. Genere par construire_v3.py.
\label{sec:predictive-robuste}
\label{sec:posture-retenue}
Une fois l'incertitude établie, il reste à choisir ce que le mémoire publie. Six postures
sont définies et chiffrées, du quantile ponctuel au pire cas à indice de queue posé. La
posture retenue est le \textbf{plug-in accompagné de sa bande}, et le motif est
réglementaire avant d'être d'opportunité : une borne haute sur un ensemble d'ambiguïté est
un suprémum sur une famille de lois, pas un quantile de la loi de perte, donc la
substituer répondrait à une autre question. Un point mérite d'être connu avant d'en
débattre : la posture robuste \textbf{est} la borne haute de l'intervalle déjà publié, de
sorte que reporter le plug-in avec sa bande publie déjà le chiffre robuste. La table des
six postures et leur bruit de simulation sont dans la version longue de ce mémoire.

\section{Un défaut de calibration, chiffré plutôt que corrigé}
% PONT v3 : la section complete est dans main_v2, version
% longue et non modifiee. Genere par construire_v3.py.
\label{sec:pu-coherence}
La calibration publiée porte une incohérence arithmétique connue : le taux de dépassement
et le nombre d'excès ne désignent pas le même échantillon, alors que la formule qui les
emploie n'a que deux degrés de liberté pour trois entrées. L'effet est \emph{calculé} et
non estimé, il vaut quelques pour cent sur le quantile unitaire, soit une fraction de la
largeur de l'intervalle publié sur cette même grandeur. La valeur reste gelée, et ce choix
est assumé : rejouer un pipeline stochastique pour un écart de cette taille déplacerait
des centaines de nombres publiés sans qu'aucun déplacement soit attribuable à la
correction, et la piste d'audit serait perdue pour un gain immatériel.

\textbf{Deux prudences à ne pas confondre.} Côté solvabilité l'écart est
anti-conservateur, il sous-estime le capital, et une sous-estimation se déclare ; côté
thèse il va dans l'autre sens, le chiffre avancé étant minoré et non gonflé. Seul le
premier engage. Le calcul, son signe et le contrôle qui le recalcule à chaque exécution
sont dans la version longue de ce mémoire.

\section{Défaillances simultanées et additivité des coûts}
\label{sec:additivite-couts}

La cascade produit, pour chaque sinistre, un \emph{ensemble} de piliers touchés, si bien que la
défaillance simultanée n'est pas un cas particulier resté de côté mais la sortie même du modèle. La
loi exacte de ce nombre est donnée au \S\ref{sec:ann-additivite} : à l'état non conforme, la
défaillance multiple est majoritaire, $62{,}31\,\%$ des sinistres touchant plus d'un
pilier contre $31{,}15\,\%$ à l'état conforme. Reste la question de l'agrégation
des coûts, et elle demande de distinguer trois énoncés que le seul mot d'additivité recouvre.

\begin{enumerate}[leftmargin=*,itemsep=2pt]
\item \textbf{Au sein d'un sinistre}, les coûts des piliers touchés \emph{s'additionnent} : c'est
une hypothèse de construction, non testée, et c'est le seul des trois que la question met en cause.
\item \textbf{En fonction des piliers non conformes}, le capital est \emph{presque additif},
$R^2=0{,}9945$ sur les $32$ configurations : c'est une mesure, et elle explique qu'une structure de
dépendance entre états ne déplace pas le capital espéré (\S\ref{sec:ann-dependance-etats}).
\item \textbf{En fonction des quatre canaux}, il est \emph{franchement super-additif},
$+5\,001$~M\euro{} d'interaction sur $14\,139$, soit $+35\,\%$ (chapitre~\ref{chap:resultats}) :
c'est également une mesure.
\end{enumerate}

\noindent Répondre \og~le modèle est super-additif~\fg{} serait donc faux : il l'est sur les canaux,
pas sur les piliers, et l'hypothèse sur les coûts d'un incident est un troisième objet.

\paragraph{Ce que la première coûte, chiffré plutôt que supposé.} On la relâche par un exposant, la
perte d'un sinistre touchant $k$ piliers valant $\bigl(\sum_j X_j\bigr)k^{\,\theta-1}$, de sorte que
$\theta=1$ redonne exactement le modèle publié et que les sinistres mono-pilier restent intouchés.
Pour $\theta$ parcourant $[0{,}70\,;1{,}30]$, l'écart entre états parcourt
$[9\,706\,;20\,487]$~M\euro, soit $76\,\%$ de l'écart publié et $15$~fois son bruit de simulation de
$\pm 734$~M\euro. L'effet est donc réel, et l'additivité des coûts est la plus lourde des
hypothèses structurelles non testées du modèle. Deux précautions pour ne pas en tirer un classement
qu'il n'établit pas : elle porte sur l'\emph{écart} là où la queue de sévérité domine l'incertitude
sur le \emph{niveau} (\S\ref{sec:trois-etages}), et c'est une hypothèse de structure et non un choix
de prudence, les trois choix posés du \S\ref{sec:pu-coherence} déplaçant davantage mais dans une
direction connue.

\paragraph{Elle n'est pas neutre entre les deux états, ce qui n'était pas attendu.} Les élasticités
valent $0{,}95$ à l'état non conforme contre $0{,}39$ à l'état conforme, un facteur $2{,}42$, parce
que l'exposant ne mord que sur les sinistres multi-piliers, majoritaires au seul état non conforme.
L'hypothèse interagit donc avec le canal de propagation au lieu de s'ajouter à côté de lui,
et atteint de ce fait la grandeur centrale du mémoire et non un simple niveau. Ce qui tient malgré
tout : sur toute l'amplitude examinée l'écart garde son signe et son ordre de grandeur, jamais
proche de zéro, si bien que la thèse ne dépend pas de l'hypothèse d'additivité et que seule son
amplitude en dépend. C'est la structure d'argument déjà employée pour les valeurs de $g$.

Le \emph{sens} de la non-additivité, lui, n'est pas déterminable : mutualisation de la remédiation
et saturation de la capacité de réponse tirent en sens contraire, et aucune source ne les sépare. Le
détail, la loi du nombre de piliers, la table complète des exposants et la raison pour laquelle
l'amplitude est posée plutôt qu'estimée sont au \S\ref{sec:ann-additivite}.

\vspace{2pt}
% SOURCES-SCRIPTS: 74 68 69


\begin{cle}
\textbf{La double lecture.} \og Insensible aux choix innocents, sensible aux leviers réels \fg :
le modèle n'est ni sur-ajusté ni fragile, il est piloté par ce qui doit le piloter.
Et la seule grande incertitude, la queue, est nommée et bornée, pas cachée. C'est ce qui
transforme \og je n'ai pas tout calibré \fg{} en \og je sais précisément de quoi mon SCR
dépend, et de quoi il ne dépend pas \fg.
\end{cle}

% =====================================================================
